%% Chapter 14 — Troubleshooting
\chapter{Troubleshooting}
\label{ch:troubleshooting}

\section{Quick Diagnostic Reference}

\begin{table}[H]
\centering
\caption{Symptom to First Action}
\begin{tabular}{L{5.5cm}L{9cm}}
\toprule
\textbf{Symptom} & \textbf{First action} \\
\midrule
OLED blank & Check 3.3V on OLED VCC; flash test sketch 04 or 05 \\
Card inserted, no reaction & Flash \texttt{06\_hc89}; confirm HC89 OUT wired to GPIO 32 \\
Card denied (should be allowed) & Check serial for denial reason; verify permissions in Firestore \\
Relay doesn't energise & Check relay LED; measure GPIO 26 with multimeter during session \\
MQTT keeps reconnecting & Check Mosquitto status on RPi; confirm MQTT\_BROKER IP \\
Sessions missing from Firestore & Check bridge logs; curl bridge health endpoint \\
Dashboard shows offline & Check node MQTT; check bridge health; check RTDB subscription in web app \\
Remote stop not working & Check RTDB /mms/commands; check bridge logs; check 5-min TTL \\
OTA stuck & Check Firebase Storage URL; check bridge OTA listener in logs \\
Node rebooting in loop & Check serial for boot reason; watchdog or OTA crash rollback \\
\bottomrule
\end{tabular}
\end{table}

\section{Serial Monitor Diagnostic Commands}

When a node is connected via USB, its Serial Monitor (115200 baud) provides the most direct diagnostic information:

\begin{table}[H]
\centering
\caption{Diagnostic Log Prefixes}
\begin{tabular}{L{2.5cm}L{11cm}}
\toprule
\textbf{Prefix} & \textbf{Meaning} \\
\midrule
\texttt{[HW]} & Hardware initialisation events (boot only) \\
\texttt{[NET]} & WiFi connection events \\
\texttt{[MQTT]} & MQTT connection, publish, subscribe, receive events \\
\texttt{[RFID]} & Card detection and read events \\
\texttt{[ACCESS]} & Access grant or denial with reason \\
\texttt{[RELAY]} & Relay state changes \\
\texttt{[SESSION]} & Session start, end, timeout events \\
\texttt{[FS]} & LittleFS operations \\
\texttt{[NVS]} & NVS reads and writes \\
\texttt{[OTA]} & OTA update progress \\
\texttt{[SYS]} & System events: boot reason, watchdog, state transitions \\
\texttt{[ERR]} & Error conditions \\
\bottomrule
\end{tabular}
\end{table}

\section{Common Issues}

\subsection{RFID Authentication Failure}

\begin{verbatim}
[RFID] Auth FAILED for UID: A1B2C3D4
\end{verbatim}

\textbf{Causes:}
\begin{itemize}[leftmargin=1.5em,noitemsep]
    \item Card written with a different master key than what is in \texttt{secrets.h}
    \item Card is not MIFARE Classic 1K (NFC tags, bank cards, transit cards fail)
    \item Card is a fresh/blank card with default Key A (\texttt{0xFFFFFFFFFFFF}) but the node tries the derived key first
\end{itemize}

\textbf{Fix}: Verify the master key is identical in all three locations. For fresh cards: the kiosk writes the derived key --- node should then authenticate with the derived key. If a card was written by a different kiosk (different master key), it cannot be read by this node.

\subsection{Schema Version Mismatch}

\begin{verbatim}
[RFID] Schema version: 0x01 (expected 0x02) — ACCESS DENIED
\end{verbatim}

\textbf{Cause}: Card was issued by V1 kiosk (old system).

\textbf{Fix}: Re-issue the card at the V2 kiosk.

\subsection{Bridge PERMISSION\_DENIED on Firestore}

\begin{verbatim}
[FIREBASE] Error writing session: 403 PERMISSION_DENIED
\end{verbatim}

\textbf{Cause}: The service account used by bridge.py does not have Firebase Admin SDK permissions.

\textbf{Fix}: Firebase Console $\rightarrow$ IAM $\rightarrow$ check the service account has ``Firebase Admin'' or ``Editor'' role. The service account email is shown in the \texttt{service-account.json} file.

\subsection{WiFi RSSI Too Low}

\begin{verbatim}
[NET] WiFi connected — IP: 192.168.0.101, RSSI: -84 dBm
\end{verbatim}

\textbf{Cause}: Node is too far from the WiFi AP.

\textbf{Impact}: Intermittent MQTT disconnects; sessions buffered in LittleFS but may not flush promptly.

\textbf{Fix}: Add a WiFi access point or range extender nearer to the machine. Acceptable RSSI: $>-80$ dBm. Ideal: $>-67$ dBm.

\section{Recovery Procedures}

See the Deployment Package (Section 16) for complete disaster recovery procedures covering:
\begin{itemize}[leftmargin=1.5em,noitemsep]
    \item Raspberry Pi hardware failure
    \item Firebase outage
    \item ESP32 hardware failure
    \item Lost master key
    \item Corrupted firmware (all nodes)
    \item Lost Firebase admin access
    \item Database corruption
\end{itemize}
